\chapter{What a Qubit Is Made Of}\label{ch:hardware}

\begin{etymology}
\textbf{System} comes from the Greek \emph{sýstēma}: \emph{syn-}, ``together'',
and \emph{histanai}, ``to set up''. Parts set up together as a whole, which we
choose to separate from the rest of the world in order to study it.

\textbf{Coherence} comes from \emph{cohaerere}: \emph{co-}, ``together'', and
\emph{haerere}, ``to stick''. In quantum physics it means that the phase
relations between amplitudes are kept, which is what allows interference.
\emph{Decoherence} is the loss of those relations, when the system gets entangled
with its surroundings and the phases leak out.
\end{etymology}

A qubit is any physical system with two levels that obeys quantum mechanics: the
spin of an electron (up or down), the polarization of a photon (horizontal or
vertical), two energy levels of an atom. Mathematically all of them are the same
vector in $\mathbb{C}^2$. Physically they could hardly be more different, and the
differences decide everything about how far a technology can go.

The hard part is not finding a two-level system. Nature is full of them. The hard
part is finding one that you can prepare, control and read, and that still keeps
its phases long enough to compute. Those demands pull against each other. A qubit
well isolated from the world keeps its coherence, but it is also hard to reach
with the controls. A qubit easy to control is, for the same reason, easy for the
environment to disturb.

This chapter is a tour of the technologies, and a checklist guides it. Every
platform is walked through the same list, item by item, so that by the end each
item means something concrete.

\section{The checklist}

\begin{history}{2000 \textperiodcentered{} IBM}
David DiVincenzo, at IBM, listed what a physical system needs to be a quantum
computer. He had written, with Daniel Loss, one of the first concrete proposals for
qubits in 1998. His list became the checklist every technology is measured against.
\end{history}

\begin{enumerate}
\item \textbf{Characterized and scalable.} You know exactly what your qubit is, and
adding more of them does not make the others worse.
\item \textbf{Initialization.} You can put every qubit in a known starting state,
$\ket{00\dots0}$.
\item \textbf{Coherence.} The qubit keeps its quantum state much longer than a gate
takes.
\item \textbf{Universal gates.} A small menu of operations is enough to build any
computation.
\item \textbf{Measurement.} You can read each qubit, reliably.
\end{enumerate}

He added two more for quantum communication: (6) converting between qubits that
stay put and qubits that travel, which in practice are photons, and (7) sending the
traveling ones faithfully from one place to another.

In 2018, John Preskill gave a name to the machines that exist today: \emph{NISQ},
noisy intermediate-scale quantum. Tens to a few thousand qubits, too noisy for error
correction at scale, too large to simulate exactly. No technology yet meets the five
criteria at the fidelity and scale that Shor's algorithm needs on a cryptographic
number. The tour shows where each one falls short.

\section{Superconducting circuits}

\begin{history}{1962--1985 \textperiodcentered{} Cambridge and Berkeley}
In 1962 Brian Josephson, a 22-year-old PhD student at Cambridge, predicted that a
current can flow without resistance through a thin insulating barrier between two
superconductors. The \emph{Josephson junction} earned him a share of the 1973
Nobel Prize. In 1984 and 1985, at Berkeley, John Clarke, Michel Devoret and John
Martinis showed that an electrical circuit built around such a junction, big enough
to see under a microscope, behaves as a single quantum object: it tunnels through
energy barriers and has discrete energy levels that they could probe with
microwaves. A circuit with billions of electrons had a quantized spectrum. They
received the Nobel Prize in Physics for it in 2025.
\end{history}

\subsection{What the qubit is}

Start with the simplest circuit there is: a capacitor and an inductor in a loop. Charge
sloshes from the capacitor through the inductor and back, and the current oscillates
at a fixed frequency. It is the mass on a spring of Chapter~\ref{ch:classical}, with
charge in place of position and current in place of velocity. Cool it until the metal
superconducts, so nothing is lost to resistance, and it becomes a quantum harmonic
oscillator: its energy can only take values on a ladder of equally spaced rungs.

\begin{reflection}{The energy of a whole circuit comes in steps?}
That sentence should bother you. Quantization was discovered in atoms. A circuit is a
piece of metal you can see under a microscope, made of billions of electrons, and we
are saying its energy cannot change smoothly.

First, which energy. Not the energy of each electron, and not heat. It is the energy
of the oscillation itself, the one that sloshes back and forth. At one instant it sits
in the electric field of the charged capacitor, $Q^2/2C$; a quarter of a period later
it sits in the magnetic field of the current through the inductor, $\Phi^2/2L$, where
$\Phi$ is the magnetic flux. Their sum is the Hamiltonian of the circuit,
\[
  H = \frac{Q^2}{2C} + \frac{\Phi^2}{2L},
\]
the same shape as the spring's $p^2/2m + kx^2/2$ in Chapter~\ref{ch:classical}: charge
plays the role of momentum, flux the role of position, and $1/C$ and $1/L$ the roles of
$1/m$ and $k$. So far this is classical physics. Three steps turn it into quantum
physics, and each one was built earlier in the book.

\textbf{Dirac's rule} (end of Chapter~\ref{ch:classical}). Replace each classical
quantity with an operator, and the Poisson bracket with a commutator. For the circuit,
$\{\Phi, Q\} = 1$ becomes
\[
  [\hat\Phi, \hat Q] = \hat\Phi\hat Q - \hat Q\hat\Phi = i\hbar .
\]
Charge and flux stop being numbers. They become operations on the state of the
circuit, and the order in which you apply them matters.

\textbf{Uncertainty} (Chapter~\ref{ch:quantum}). Two operators that do not commute
have no eigenvectors in common, so no state of the circuit has a definite charge and a
definite flux at the same time. Robertson's inequality gives
$\Delta\Phi\,\Delta Q \geq \hbar/2$. The circuit can never be perfectly at rest, with
no charge and no current at once.

\textbf{Energies are eigenvalues} (Chapter~\ref{ch:quantum}, third postulate). The
energy is now the operator $\hat H$, and a measurement of it can only return one of
its eigenvalues. For the particle in a box they came out discrete because only waves
that fit between the walls are allowed. The circuit is a box with soft walls: the
energy $\Phi^2/2L$ grows on both sides, so the wave describing the state has to die
out on both sides, and only some shapes manage it. Solving for those shapes gives a
ladder of energies, equally spaced, $\hbar\omega$ apart, with
$\omega = 1/\sqrt{LC}$.

If one of these steps lost you, the problem is not in this box. You skipped
Chapter~\ref{ch:quantum}, or the end of Chapter~\ref{ch:classical}, where each step is
built from its history. Go back; this box only applies them to a circuit.

So what is a rung of the ladder? Not a place, and not a part of the circuit. Each rung
is one value that the total energy of the oscillation is allowed to have:
\[
  E_n = \hbar\omega\left(n + \tfrac12\right), \qquad n = 0, 1, 2, \dots
\]
The number $n$ counts how many quanta of oscillation the circuit is holding, and each
quantum carries exactly the energy of one microwave photon of frequency $\omega/2\pi$.
Climbing one rung means absorbing one such photon; dropping one means emitting it. There
is nothing in between: the circuit cannot hold half a quantum. Even the bottom rung is not
zero. The circuit keeps $\tfrac12\hbar\omega$ there, the \emph{zero-point energy}, because
charge and flux cannot both be exactly zero at once; that would violate the uncertainty
principle of Chapter~\ref{ch:quantum}.

Second, why nobody sees the steps in ordinary electronics. They are there, but tiny.
At 5~GHz one step is $3.3\times10^{-24}$~J. A capacitor of 1~pF charged to 1~V stores
$5\times10^{-13}$~J, which is 150 billion steps: a ramp made of grains of sand looks
smooth. At room temperature, thermal agitation is 1250 times larger than one step, so
the state is smeared over a wide band of rungs and no single one can be picked out. And
resistance drains the oscillation before any step could be noticed.

\begin{center}
\begin{tikzpicture}[x=1mm, y=1mm]
  % ---------- ordinary circuit, warm ----------
  \begin{scope}
    \foreach \i in {0,...,64} \draw[black!35, line width=0.25pt] (0,\i*0.5) -- (22,\i*0.5);
    \shade[top color=hfAlert!0, bottom color=hfAlert!0, middle color=hfAlert!55, opacity=0.75]
      (-0.5,9) rectangle (22.5,27);
    \draw[black!60, decorate, decoration={brace, amplitude=3pt}] (23.5,27) -- (23.5,9)
      node[midway, right=3pt, hf/note, align=left] {thermal\\spread:\\$\sim$1250\\rungs};
    \node[font=\sffamily\small, above, align=center] at (11,33.5) {ordinary circuit, 300~K};
    \node[hf/note, below, align=center] at (11,-1.5) {rungs too close to tell apart;\\the state is smeared\\(not to scale)};
  \end{scope}
  % arrow
  \draw[black!40, -{Stealth[length=2.4mm]}, line width=0.8pt] (44,16) -- (70,16)
    node[midway, above, hf/note, align=center] {superconducting,\\cooled to 15~mK};
  % ---------- same circuit, cold ----------
  \begin{scope}[xshift=92mm]
    \foreach \n/\e in {0/{\tfrac12}, 1/{\tfrac32}, 2/{\tfrac52}, 3/{\tfrac72}} {
      \draw[black!70, line width=1pt] (0,\n*9) -- (20,\n*9);
      \node[font=\small, left] at (0,\n*9) {$n = \n$};
      \node[font=\small, right] at (20,\n*9) {$\e\,\hbar\omega$};
    }
    \fill[hfConcept] (10,1.6) circle (1.6);
    \draw[hfConcept, -{Stealth[length=1.8mm]}, line width=0.8pt] (15,0.6) -- (15,8.4);
    \node[hf/note, text=hfConcept, anchor=west] at (15.5,4.5) {one photon};
    \node[font=\sffamily\small, above, align=center] at (10,33.5) {same circuit, cold};
    \node[hf/note, below, align=center] at (10,-1.5) {each rung is one allowed energy;\\the circuit sits on the bottom one};
  \end{scope}
\end{tikzpicture}
\end{center}

The transmon removes all three disguises. It is cooled far below the size of one step,
it is superconducting so nothing drains away, and it holds zero or one quantum, not
billions. What remains is the bare ladder.

The last reason it works is superconductivity itself. Below the critical temperature,
the electrons pair up and all the pairs share a single quantum state. Billions of
particles move as one, so the whole circuit has a single degree of freedom, and that
single degree of freedom behaves like one particle on a spring. In 1980 Anthony Leggett
asked whether quantum mechanics would still hold for a variable like that, describing a
macroscopic object. The experiments of Clarke, Devoret and Martinis in 1985 were the
answer: it does.
\end{reflection}

That ladder is useless as a qubit. A microwave pulse tuned to the jump from the bottom
rung to the next has exactly the frequency of every other jump too, so it pushes the
circuit up the ladder and out of any pair of levels you wanted to use.

Now replace the inductor with a Josephson junction. The junction behaves like an
inductor whose value depends on the current through it, a \emph{nonlinear} one. It
bends the oscillator so that the rungs get closer together as you climb. The first jump
and the second now have different frequencies, and a pulse tuned to the first leaves
the second alone.

\begin{figure}[htbp]
\centering
\begin{tikzpicture}[x=1mm, y=1mm]
  % ---------- (a) LC oscillator ----------
  \begin{scope}
    \draw[line width=0.8pt] (0,4) -- (0,0) -- (16,0) -- (16,3);
    \draw[line width=0.8pt] (0,6) -- (0,10) -- (16,10) -- (16,7);
    \draw[line width=1.2pt] (-3,4) -- (3,4) (-3,6) -- (3,6);
    \draw[line width=0.8pt, decorate, decoration={coil, aspect=0.6, segment length=1.6mm, amplitude=1.6mm}] (16,7) -- (16,3);
    \node[hf/note, left] at (-3,5) {$C$};
    \node[hf/note, right] at (18,5) {$L$};
    \node[font=\sffamily\small, above] at (8,11) {(a) LC oscillator};
    % ladder
    \foreach \lvl/\name in {0/0, 9/1, 18/2, 27/3} {
      \draw[black!70, line width=1pt] (0,-40+\lvl) -- (16,-40+\lvl);
      \node[font=\small, left] at (0,-40+\lvl) {$\ket{\name}$};
    }
    \foreach \a/\b in {0/9, 9/18, 18/27}
      \draw[hfConcept, -{Stealth[length=1.8mm]}, line width=0.8pt] (8,-40+\a+0.6) -- (8,-40+\b-0.6);
    \node[hf/note, below, align=center] at (8,-42) {every jump the same:\\a pulse climbs the whole ladder};
  \end{scope}
  % arrow between panels
  \draw[black!40, -{Stealth[length=2.4mm]}, line width=0.8pt] (27,5) -- (51,5)
    node[midway, above, hf/note] {swap $L$ for a junction};
  % ---------- (b) transmon ----------
  \begin{scope}[xshift=70mm]
    \draw[line width=0.8pt] (0,4) -- (0,0) -- (16,0) -- (16,10) -- (0,10) -- (0,6);
    \draw[line width=1.4pt] (-5,4) -- (5,4) (-5,6) -- (5,6);
    \draw[hfAlert, line width=1.2pt] (13.5,3.5) -- (18.5,6.5) (13.5,6.5) -- (18.5,3.5);
    \node[hf/note, left] at (-5,5) {large $C$};
    \node[hf/note, right, text=hfAlert] at (19,5) {junction};
    \node[font=\sffamily\small, above] at (8,11) {(b) transmon};
    % ladder
    \fill[hfConcept!10] (-1,-41.5) rectangle (17,-30);
    \foreach \lvl/\name in {0/0, 9/1, 16.5/2, 22.5/3} {
      \draw[black!70, line width=1pt] (0,-40+\lvl) -- (16,-40+\lvl);
      \node[font=\small, left] at (0,-40+\lvl) {$\ket{\name}$};
    }
    \draw[hfConcept, -{Stealth[length=1.8mm]}, line width=1pt] (8,-39.4) -- (8,-31.6);
    \draw[hfAlert, -{Stealth[length=1.8mm]}, line width=0.8pt, dashed] (8,-30.4) -- (8,-24.1);
    \draw[black!30, -{Stealth[length=1.8mm]}, line width=0.8pt] (8,-22.9) -- (8,-18.1);
    \draw[hfConcept, line width=0.8pt, decorate, decoration={brace, amplitude=4pt, mirror}]
      (18.5,-40) -- (18.5,-31);
    \node[hf/note, anchor=west, text=hfConcept, align=left] at (22,-35.5) {\textbf{the qubit}\\$\approx 5$ GHz};
    \node[hf/note, anchor=west, text=hfAlert, align=left] at (22,-27) {$\ket1\to\ket2$ is 200--300 MHz\\lower: the pulse misses it};
    \node[hf/note, below, align=center] at (8,-42) {each jump a bit smaller};
  \end{scope}
\end{tikzpicture}
\caption{What a superconducting qubit is. (a) A capacitor and an inductor make an
oscillator whose energy levels are equally spaced. (b) Replace the inductor with a
Josephson junction and the levels become uneven. The qubit is the two lowest levels of
this circuit.}\label{fig:ladder}
\end{figure}

So the qubit is not a particle or a little object on the chip. It is the two lowest
energy states of a whole electrical circuit, a few hundred micrometers across, with
billions of paired electrons moving as one. $\ket0$ is the circuit at rest in its
ground state. $\ket1$ is the circuit holding one quantum of oscillation, an energy of
about 20 micro-electronvolts, which is what a single microwave photon at 5~GHz
carries. A superposition $\alpha\ket0 + \beta\ket1$ is the circuit in both at once.

The first circuits of this kind came in quick succession. In 1999, at NEC in Tsukuba,
Yasunobu Nakamura, Yuri Pashkin and Jaw-Shen Tsai showed coherent oscillations in a
``Cooper-pair box'', a tiny island of superconductor; they lasted a few nanoseconds.
The same year Hans Mooij's group in Delft proposed the flux qubit, a loop with current
flowing both ways at once, seen in experiments in 2000. In 2002, at Saclay, Denis Vion,
Daniel Esteve, Michel Devoret and colleagues built the \emph{quantronium}, and at NIST
in Boulder, Martinis's group demonstrated the phase qubit. In 2007, at Yale, Jens Koch,
Robert Schoelkopf, Steven Girvin and colleagues proposed the \emph{transmon}: a
Cooper-pair box with a large capacitor added, which makes it nearly immune to stray
charges at the price of a smaller nonlinearity. The transmon became the dominant
design, used by IBM and Google among others. Now the checklist.

\subsection*{1. Characterized and scalable}

\emph{Characterized} means you know the Hamiltonian of your qubit: the energies of its
levels, how it couples to its neighbors, and what else is nearby. For a transmon that is
the ladder of Figure~\ref{fig:ladder}. The qubit frequency is typically around 5~GHz,
the next jump is 200 to 300~MHz lower, and the third rung is always there, waiting for
a pulse that is too strong or too short. A transmon that ends up in $\ket2$ has left the
computation, an error called \emph{leakage}.

Transmons are made, not found, and no two come out of fabrication identical. Each one
has its own frequency and must be measured and calibrated on its own, and the
calibration drifts. Characterization is not done once; it is a daily job.

\emph{Scalable} means adding qubits does not make the rest worse. The transmon's best
argument is here: chips are made with the lithography of ordinary electronics, so
drawing a thousand qubits costs little more than drawing ten. Its worst problem is also
here. Every qubit needs its own control lines, running from room temperature down into
the refrigerator, and each wire brings heat and noise with it. Hundreds of qubits fit;
a million wires do not.

\subsection*{2. Initialization}

Before a computation, every qubit has to be in $\ket0$. Nature does that for free if
the qubit is cold enough, because a system in equilibrium sits in its lowest level
except for thermal kicks. The kicks are small when the temperature is small compared
with the energy of the jump. For a 5~GHz transmon, that energy corresponds to a
temperature of 0.24~K. Boltzmann's law gives the chance of finding the qubit in $\ket1$:

\begin{center}
\begin{tabular}{lr}
\hline
temperature & qubit found in $\ket1$ \\
\hline
room temperature, 300~K & 49.98\% \\
liquid helium, 4~K & 48.5\% \\
50~mK & 0.8\% \\
15~mK & 0.00001\% \\
\hline
\end{tabular}
\end{center}

At room temperature the qubit is a coin toss. At 4~K, the temperature of liquid helium,
it is still a coin toss. Only well below 0.24~K does it settle in $\ket0$, and that is
why superconducting chips live in dilution refrigerators at 10 to 20~millikelvin, colder
than outer space. In practice the chip runs warmer than the refrigerator, and a fraction
of a percent of qubits start in $\ket1$, so machines often reset them actively before
each run: measure, and flip the ones that came out wrong.

\subsection*{3. Coherence}

A transmon in $\ket1$ eventually drops to $\ket0$, handing its quantum of energy to the
environment. The probability that it has survived after a time $t$ is $e^{-t/T_1}$, where
$T_1$ is the \emph{relaxation time}. The relative phase between $\ket0$ and $\ket1$ also
drifts and is lost, on a time scale called $T_2$. Interference needs both.

What the criterion asks for is a ratio, not a time. A transmon with $T_1 = 100\,\mu$s
running 50~ns gates loses about 0.05\% per gate to relaxation alone, and fits about 2000
gates into one $T_1$. The best transmons reported in 2025, with $T_1$ around a
millisecond, fit about 20\,000. In 1999 Nakamura's box lasted a few nanoseconds; in
twenty-five years coherence improved by about six orders of magnitude.

Shor's algorithm on RSA-2048 needs billions of operations. A ratio of $10^4$ is not
close, and no improvement in materials will close five orders of magnitude. That is why
error correction, the subject of Part~\ref{part:qec}, is not optional.

\subsection*{4. Universal gates}

A one-qubit gate on a transmon is a microwave pulse at the qubit frequency. While the
pulse is on, the state turns steadily from $\ket0$ toward $\ket1$ and back, a rotation on
the Bloch sphere; the length of the pulse sets the angle and its phase sets the axis.

\begin{history}{1937 \textperiodcentered{} Columbia}
Isidor Rabi sent beams of molecules through an oscillating magnetic field and saw their
nuclear spins flip when the field was tuned to the right frequency. The steady flipping
back and forth under a resonant drive is called a \emph{Rabi oscillation}, and Rabi's
method won him the 1944 Nobel Prize. It became nuclear magnetic resonance, then MRI.
Nakamura's result of 1999 was a Rabi oscillation in a circuit, and every single-qubit gate
on every platform since is one too.
\end{history}

Two-qubit gates need two transmons to interact, through a capacitor or a tunable coupler
between them. That only works between neighbors on the chip. A gate between distant
qubits has to be built from a chain of swaps, which adds time and error.

Is that menu enough? \emph{Universal} means that a small, fixed set of gates can
approximate any unitary to any precision, and the right set is not obvious.

\begin{history}{1998 \textperiodcentered{} the Gottesman--Knill theorem}
Daniel Gottesman, building on work with Emanuel Knill, proved that circuits made only of
Hadamard, CNOT and the phase gate $S = \mathrm{diag}(1, i)$, plus measurements, can be
simulated efficiently on a classical computer, no matter how many qubits they use or how
much entanglement they create. These are the \emph{Clifford} gates, and with them alone a
quantum computer is no more powerful than a laptop.
\end{history}

Add a single gate, $T = \mathrm{diag}(1, e^{i\pi/4})$, a $45^\circ$ turn about the vertical
axis of the Bloch sphere, and the set becomes universal. Robert Solovay and Alexei Kitaev
showed independently, in 1995 and 1997, that any one-qubit gate can then be approximated
to precision $\varepsilon$ by a sequence of $H$ and $T$ gates whose length grows only like
a power of $\log(1/\varepsilon)$. A transmon can do any rotation directly, so it passes this
criterion easily today. Under error correction it will not: rotations become sequences of
discrete gates, and the $T$ gate is by far the most expensive of them.

\subsection*{5. Measurement}

A transmon is read through a microwave resonator coupled to it. The resonator's frequency
shifts slightly depending on whether the qubit is in $\ket0$ or $\ket1$, so a weak pulse
bounced off it comes back changed. This scheme, \emph{circuit quantum electrodynamics},
was proposed by Alexandre Blais and colleagues at Yale in 2004. It takes hundreds of
nanoseconds, during which the qubit may relax and give the wrong answer.

Readout errors compound. If each qubit is read correctly 99\% of the time, all 100 qubits
of a register come out right only 37\% of the time. At 99.9\% the chance rises to 90\%.
Every fraction of a percent at the end of a computation is paid again for every qubit.

\subsection*{6 and 7. Moving qubits}

Here the transmon is weakest. Its quantum lives in microwave photons, which cannot leave
the refrigerator without being swamped by thermal noise; at room temperature, as the
table above shows, a microwave mode is a coin toss. Linking two refrigerators needs a
converter from microwave to optical photons that preserves quantum states, and in 2025
that was still a laboratory problem.

\section{Trapped ions}

\begin{history}{1995 \textperiodcentered{} Innsbruck and Boulder}
Ignacio Cirac and Peter Zoller, in Innsbruck, proposed the first concrete architecture:
ions held in a line by electric fields in a vacuum, with two levels of each ion as the
qubit, and laser pulses as the gates. The ions interact through their shared vibrations.
The same year, Christopher Monroe and David Wineland, at NIST in Boulder, demonstrated a
controlled-NOT gate on a single beryllium ion. Wineland shared the 2012 Nobel Prize, with
Serge Haroche, for manipulating individual quantum systems.
\end{history}

Here the qubit is an atom that has lost one electron. Its outermost electron has two
energy levels that differ by the orientation of its spin relative to the nucleus. In the
ytterbium ion $^{171}\mathrm{Yb}^+$, one of the most used, the two levels are 12.6~GHz
apart. Nature built the ladder; nobody had to bend it. The same checklist gives a very
different picture.

\subsection*{1. Characterized and scalable}

Every ytterbium ion in the universe is identical, so every qubit has exactly the same
frequency, known to many decimal places. Characterization comes for free, and this is the
ions' great advantage over anything fabricated.

Scaling is their problem. The ions sit in a line, held apart by their mutual repulsion and
coupled by the vibrations of the whole chain. Add ions and the chain gets longer, its
vibrations crowd together, and gates get slower and harder to aim. A single chain stops
being practical at a few dozen ions. The plan since 2002, proposed by David Kielpinski,
Monroe and Wineland, is to shuttle ions between zones of a trap, like trains in a yard, so
that gates always act on small groups.

\subsection*{2. Initialization}

Compare with the transmon. A 12.6~GHz qubit at room temperature is in $\ket1$ about half the
time, exactly the coin toss of the table above, and ion traps run at room temperature. They
do not need cold, because they do not rely on equilibrium. A laser tuned to excite the ion
from every state except $\ket0$ keeps kicking it until it falls into $\ket0$, where the laser
can no longer touch it. This \emph{optical pumping} prepares the qubit with probability above
99.9\%. The criterion asks for a known state, not a cold one; a one-way process is enough.

\subsection*{3. Coherence}

The two levels of the ion are barely disturbed by stray fields, and the ion floats in vacuum,
touching nothing. In 2021 a single ion kept its phase for over an hour. That is about a
billion times longer than a transmon.

But criterion 3 is about a ratio. A two-qubit gate on ions takes around 100~$\mu$s, a
thousand times slower than on a transmon, so the number of gates per coherence time lands at
about $10^4$, the same range. The hour-long coherence does not buy a longer computation; it
buys a slower one.

\subsection*{4. Universal gates}

One-qubit gates are laser or microwave pulses, Rabi oscillations again. Two-qubit gates use
the chain: a laser couples the internal state of two ions to the vibration they share, and
the vibration carries the interaction from one to the other. Any ion in the chain can be
coupled to any other, so there are no swaps. In 2024, Quantinuum reported two-qubit gates
with 99.9\% fidelity, the best of any platform at the time.

\subsection*{5. Measurement}

Ions are read by a method Hans Dehmelt proposed in 1975. Shine a laser tuned so that the ion
scatters light only when it is in one of the two states. In that state it glows, millions of
photons per second; in the other it stays dark. Count the photons and you have the answer,
with fidelity above 99.9\%.

\subsection*{6 and 7. Moving qubits}

Ions emit optical photons, which travel through ordinary fiber. In 2007 Monroe's group
entangled two ions held in separate traps a meter apart, using the photons each one
emitted. This is the ions' second advantage: modular machines linked by light are a natural
extension, not a research problem.

\section{The others}\label{sec:others}

The other platforms are easiest to understand through the criterion where they break. Topological
qubits come after them, in a section of their own, because they need more explanation.

\textbf{Nuclear magnetic resonance: criteria 1 and 2.} In 1997 David Cory, Amr Fahmy and
Timothy Havel, and separately Neil Gershenfeld and Isaac Chuang, proposed using the nuclear
spins of molecules in a liquid as qubits, driven by the radio pulses of NMR machines. In 2001
Lieven Vandersypen and colleagues at IBM used seven spins to run Shor's algorithm and factor
15. But a nuclear spin at 400~MHz, at room temperature, is in its upper state with probability
49.998\%: initialization fails, and badly. NMR worked around it with statistical tricks over
trillions of molecules, at the cost of a signal that falls exponentially with the number of
qubits, which fails criterion 1. In 1999 Samuel Braunstein and colleagues showed that the NMR
experiments of the time never created entanglement at all. It was useful for testing ideas, and
a dead end for computing.

\textbf{Neutral atoms: criteria 1 and 4.} Atoms without charge, held by tightly focused laser
beams called optical tweezers. Like ions, they are identical by nature, and they can be moved,
which gives flexible connections. Scale is their strength: in 2025 a Caltech group trapped more
than six thousand atoms in one array. Trapping is not computing, though. Neutral atoms barely
interact, so two-qubit gates excite them to huge Rydberg states, an idea proposed by Dieter
Jaksch and colleagues in 2000, and those gates are slower and less accurate than one-qubit gates.
Atoms also escape the tweezers, which is a qubit vanishing from the computation.

\textbf{Photons: criterion 4.} The qubit is the polarization or the path of a photon. Photons
barely interact with the environment, which is excellent for criterion 3 and makes them the
natural answer to criteria 6 and 7. For the same reason they barely interact with each other,
which is fatal for two-qubit gates. In 2001 Emanuel Knill, Raymond Laflamme and Gerard Milburn
showed that gates can be built from mirrors, beam splitters and measurements, but they only work
with some probability, and the overhead is large.

\textbf{Spins in semiconductors: criterion 1.} Daniel Loss and DiVincenzo proposed in 1998 to use
the spin of single electrons trapped in quantum dots, tiny traps in silicon. They would scale like
transistors, which is the whole appeal. But each dot is shaped by a few atoms of disorder, so every
qubit is slightly different, and characterizing and controlling them uniformly is still the
obstacle.

\section{Topological qubits}\label{sec:topological}

Every platform so far stores the qubit in something local: two levels of a circuit, of an ion,
of an atom, of a spin. Then it defends that local thing from the world, first by isolating it
and in the end by error correction, which Part~\ref{part:qec} will show costs hundreds of
physical qubits per logical one. The topological approach tries to store the qubit somewhere
local noise cannot reach. It has not yet passed criterion 1: there is no undisputed
demonstration that the qubit exists, and Chapter~\ref{ch:majorana} tells that story. This
section explains how the device is supposed to work, because the claims in that chapter only
make sense against it. It uses the bosons, fermions and anyons of Section~\ref{sec:identical}.

\begin{concept}{Topological protection}
Information stored in a global property of a system, such as the joint parity of two
distant zero modes, so that no local disturbance can change it.

\textbf{What it fixes:} criterion 3 by construction. Instead of fighting decoherence only
with many physical qubits per logical qubit, the hardware itself does part of the job.

\textbf{What it costs:} materials that nobody knew how to make, a protection that only
holds while an energy gap stays open and no electron jumps in from outside, and a gate set
that is not universal.
\end{concept}

\subsection*{The gap}

In a superconductor the electrons bind into Cooper pairs, and the pairs, being bosons, all
share one collective state. Leon Cooper showed in 1956 that the pairs form, and in 1957 John
Bardeen, Cooper and Robert Schrieffer built the full theory. Breaking a pair costs energy: at
least $\Delta$ for each electron set free, the \emph{superconducting gap}. Below the gap there
is nothing to excite, which is why a superconducting circuit is so quiet. The BCS theory gives
$\Delta \approx 1.76\,k_B T_c$, with $T_c$ the critical temperature. Aluminum, with $T_c =
1.2$~K, has a gap of about $180\,\mu$eV. Lead, with $T_c = 7.2$~K, has about 1.4~meV,
somewhat more than the BCS estimate of 1.1~meV.

At 15~mK, $k_B T$ is $1.3\,\mu$eV, a hundred and forty times smaller than the gap of
aluminum, so heat alone should break essentially no pairs. In practice a few pairs are always
broken anyway, by stray radiation and by heat leaking down the wiring, far more than the
temperature predicts. The free electrons they leave, called \emph{quasiparticles}, wander
through the metal. For a transmon they are one source of error among many. For a topological
qubit they are the main enemy, as the next figure shows.

\subsection*{Half of an electron}

An excitation of a superconductor is not quite an electron. It is a mixture of an electron and
a \emph{hole}, the absence of an electron, which behaves as a particle of opposite charge. A
hole is to an electron what an antiparticle is to a particle. In some superconductors a special
excitation can appear that is an exact half-and-half mixture, at zero energy. It is its own
antiparticle, which is the property Ettore Majorana found in an equation in 1937, and so it is
called a \emph{Majorana zero mode}. Its operator equals its own adjoint, $\gamma = \gamma^\dagger$.

One Majorana zero mode is half of an ordinary fermion state. Two of them, $\gamma_1$ and
$\gamma_2$, combine into one complete fermion mode,
\[
  c = \frac{\gamma_1 + i\gamma_2}{2},
\]
and an ordinary fermion mode can only be empty or full, by the exclusion principle. That
``empty or full'' is a bit, and its value is called the \emph{parity}.

The point is where the bit lives. It is not in either half. It is in the pair, and the halves
can be micrometers apart. Noise that touches one end of the wire, a stray field, a vibrating
atom, cannot read or change a value that only exists in the combination of both ends. One
thing can change it: an electron arriving from outside, such as a quasiparticle, which turns
an empty mode into a full one. This is called \emph{quasiparticle poisoning}, and the average
wait for it is the parity lifetime that Chapter~\ref{ch:majorana} is about.

\begin{figure}[htbp]
\centering
\begin{tikzpicture}[x=1mm, y=1mm]
  % where each half lives
  \draw[black!25, dashed, line width=0.6pt] (0,13) -- (90,13);
  \fill[hfConcept!15] (0,13) -- plot[domain=0:40, samples=60] (\x, {13 + 8*exp(-\x/5)}) -- (40,13) -- cycle;
  \fill[hfConcept!15] (50,13) -- plot[domain=50:90, samples=60] (\x, {13 + 8*exp(-(90-\x)/5)}) -- (90,13) -- cycle;
  \draw[hfConcept, line width=1pt, domain=0:40, samples=60] plot (\x, {13 + 8*exp(-\x/5)});
  \draw[hfConcept, line width=1pt, domain=50:90, samples=60] plot (\x, {13 + 8*exp(-(90-\x)/5)});
  \node[hf/note] at (45,22) {each half sits at one end;\\the middle holds nothing};
  % the device
  \draw[hfSource, fill=black!30, line width=0.8pt] (0,6) rectangle (90,10);
  \node[font=\sffamily\footnotesize, text=white] at (45,8) {superconductor};
  \draw[hfSource, fill=hfSource!10, line width=0.8pt] (0,0) rectangle (90,6);
  \node[font=\sffamily\footnotesize, text=black!70] at (45,3) {semiconductor nanowire};
  \draw[black!50, -{Stealth[length=2mm]}, line width=0.8pt] (52,30) -- (78,30);
  \node[hf/note, anchor=east] at (51,30) {magnetic field along the wire};
  % the two halves
  \fill[hfConcept] (5,3) circle (2.4);
  \fill[hfConcept] (85,3) circle (2.4);
  \node[font=\small, text=white] at (5,3) {$\gamma_1$};
  \node[font=\small, text=white] at (85,3) {$\gamma_2$};
  % together: one mode
  \draw[hfConcept, dashed, line width=0.8pt] (5,0.6) .. controls (30,-18) and (60,-18) .. (85,0.6);
  \node[hf/note, text=hfConcept] at (45,-17.5) {together, one fermion mode $c = (\gamma_1 + i\gamma_2)/2$:\\empty or full, and that is the parity bit};
  % local noise
  \draw[hfAlert, line width=0.8pt, decorate, decoration={snake, amplitude=1mm, segment length=3mm}] (-17,3) -- (-4,3);
  \draw[hfAlert, -{Stealth[length=2mm]}, line width=0.8pt] (-4,3) -- (2.2,3);
  \node[hf/note, text=hfAlert, anchor=north] at (-11,-1) {local noise\\touches one half:\\the bit is safe};
  % poisoning
  \draw[hfAlert, -{Stealth[length=2mm]}, line width=0.8pt] (101,3) -- (87.8,3);
  \node[draw=hfAlert, circle, inner sep=0.6pt, font=\scriptsize, text=hfAlert] at (104,3) {$e^-$};
  \node[hf/note, text=hfAlert, anchor=north] at (101,-1) {an electron\\from outside\\flips the bit};
\end{tikzpicture}
\caption{A bit stored in the two ends of a wire. Two Majorana zero modes, $\gamma_1$ and
$\gamma_2$, sit at the ends of a semiconductor nanowire covered by a superconductor, in a
magnetic field along the wire. Neither is a state by itself; together they form one fermion
mode, empty or full, and that is the parity bit. Noise acting on one end reaches only half of
it and cannot read or change the bit. An electron arriving from outside does change it, and
the average wait for that event is the parity lifetime.}\label{fig:majorana-wire}
\end{figure}

\subsection*{The recipe}

In October 2000 Kitaev wrote down the simplest model with a Majorana zero mode at each end: a
chain of sites in which electrons of a single spin direction pair up. Ordinary superconductors
pair electrons of opposite spins, so no known material did that. In 2010 two groups found a way
to fake it with three ordinary ingredients: Roman Lutchyn, Jay Sau and Sankar Das Sarma, and
separately Yuval Oreg, Gil Refael and Felix von Oppen.

\begin{itemize}
\item A semiconductor nanowire with strong \emph{spin--orbit coupling}, such as indium arsenide,
in which an electron's spin is tied to its direction of motion.
\item A magnetic field along the wire, which pushes one spin direction up in energy and out of
the way.
\item A coat of ordinary superconductor, such as aluminum, which by contact induces pairing in
the wire. This is the \emph{proximity effect}.
\end{itemize}

Past a critical magnetic field, the wire enters a \emph{topological phase}, and a Majorana zero
mode should appear at each end. The phase has its own gap, the \emph{topological gap}, which
protects the modes, and it is much smaller than the gap of the superconductor that induces it:
tens of microelectronvolts in Microsoft's devices, against $180\,\mu$eV for aluminum. Two
things go wrong easily. Disorder in the wire can close the topological gap in places, and then
ordinary states trapped near the ends, called \emph{Andreev bound states}, can sit near zero
energy and imitate Majorana modes in most measurements. And the two modes of a finite wire are
not perfectly independent: they overlap a little, by an amount that falls exponentially with
the length of the wire, and that overlap slowly scrambles the phase of the qubit.

\subsection*{The tetron}

One wire gives one parity, which is a bit, not a qubit. In 2017 Torsten Karzig and colleagues at
Microsoft published designs for qubits made of four or more modes, and the simplest is the one
Microsoft builds: the \emph{tetron}. Two parallel topological wires are joined by an ordinary
superconducting backbone into a single small island, with four modes, $\gamma_1$ to $\gamma_4$,
at the four ends.

The island is small, so adding one electron to it costs a noticeable amount of electrostatic
energy, the \emph{charging energy}. That fixes the total parity of the four modes and is the first
defense against poisoning. With the total fixed, the four modes have two states left, and those
are one qubit. Its $Z$ is the parity of the pair along one wire, $i\gamma_1\gamma_2$. Its $X$ is
the parity of a pair across the two wires, $i\gamma_1\gamma_3$. A line of algebra with
$\gamma^2 = 1$ and $\gamma_j\gamma_k = -\gamma_k\gamma_j$ shows that these two operators
anticommute, exactly like the $X$ and $Z$ of Chapter~\ref{ch:state}.

Reading a parity uses a \emph{quantum dot}, a tiny island in the semiconductor that holds a few
electrons, coupled to the two modes being measured. An electron can tunnel from the dot through
either mode, and the two paths close a loop threaded by a magnetic flux, so they interfere. How
easily charge moves on and off the dot, its \emph{quantum capacitance}, then depends on the
parity of the pair, and a radio-frequency signal reflected off the dot reads it in a single shot.
A qubit needs both loops: one along a wire for $Z$ and one across the wires for $X$.

\begin{figure}[htbp]
\centering
\begin{tikzpicture}[x=1mm, y=1mm]
  % wires
  \draw[hfSource, fill=hfSource!10, line width=0.8pt] (0,20) rectangle (70,25);
  \draw[hfSource, fill=hfSource!10, line width=0.8pt] (0,0) rectangle (70,5);
  \draw[hfSource, fill=black!30, line width=0.8pt] (32,5) rectangle (38,20);
  \node[hf/note, anchor=west] at (39,12.5) {superconducting backbone};
  \node[font=\sffamily\footnotesize, text=black!70] at (50,22.5) {topological wire};
  \node[font=\sffamily\footnotesize, text=black!70] at (50,2.5) {topological wire};
  % modes
  \foreach \x/\y/\n in {4/22.5/1, 66/22.5/2, 4/2.5/3, 66/2.5/4} {
    \fill[hfConcept] (\x,\y) circle (2.3);
    \node[font=\scriptsize, text=white] at (\x,\y) {$\gamma_{\n}$};
  }
  % Z loop along the top wire
  \draw[hfConcept, dashed, line width=0.9pt, rounded corners=4pt] (-3,18) rectangle (73,33);
  \fill[hfConcept!25] (35,33) circle (2.6);
  \draw[hfConcept, line width=0.8pt] (35,33) circle (2.6);
  \node[font=\scriptsize] at (35,33) {QD};
  \node[hf/note, text=hfConcept, anchor=west] at (75,29) {$Z$ loop: a dot reads\\the parity $i\gamma_1\gamma_2$\\along one wire};
  % X loop across the left ends
  \draw[hfReflect, dashed, line width=0.9pt, rounded corners=4pt] (-6,-3) rectangle (11,28.5);
  \fill[hfReflect!25] (-6,12.5) circle (2.6);
  \draw[hfReflect, line width=0.8pt] (-6,12.5) circle (2.6);
  \node[font=\scriptsize] at (-6,12.5) {QD};
  \node[hf/note, text=hfReflect, anchor=east] at (-10,12.5) {$X$ loop: a dot reads\\the parity $i\gamma_1\gamma_3$\\across the wires};
  % note
  \node[hf/note] at (35,-9) {one island, four modes, total parity fixed by the charging energy: one qubit};
\end{tikzpicture}
\caption{A tetron, the qubit Microsoft builds. Two topological wires, joined by an ordinary
superconducting backbone, hold four Majorana zero modes. With the total parity fixed, the
parity along one wire is the qubit's $Z$ and the parity across the wires is its $X$. Each is
read by a quantum dot (QD) coupled to the two modes, through an interference loop threaded by a
magnetic flux.}\label{fig:tetron}
\end{figure}

\subsection*{Gates by measurement}

Majorana zero modes are non-abelian anyons. Dmitri Ivanov showed in 2000 that exchanging two of
them applies a unitary on the shared states, not just a phase, which is what Kitaev's 1997
proposal needed. Physically moving modes along wires is hard, though, and in 2008 Parsa
Bonderson, Michael Freedman and Chetan Nayak showed that the same gates can be produced by a
sequence of parity measurements, with no motion at all. That is Microsoft's plan:
\emph{measurement-only} topological computing, with the dots of Figure~\ref{fig:tetron} as the
only tool.

The limit is the gate set. Exchanging Majorana modes produces only Clifford gates, quarter turns,
and the Gottesman--Knill theorem earlier in this chapter says that Clifford gates alone can be
simulated on a laptop. The $T$ gate has to come from \emph{magic states}, special states prepared
outside the topological protection and purified at a cost that error correction has to pay.
Topological protection buys criterion 3 and a smaller overhead per logical qubit. It does not
buy universality.

\subsection*{The checklist}

\textbf{1. Characterized and scalable.} Not shown. Whether the modes in the devices built so far
are topological or ordinary Andreev states is the open question of Chapter~\ref{ch:majorana}. If
they are topological, the wires are made by lithography, like transmons, and could scale the same
way.

\textbf{2. Initialization.} By measurement: measuring a parity leaves the qubit in a known state.

\textbf{3. Coherence.} Two clocks, one per parity. $Z$ survives until a quasiparticle poisons the
island. $X$ survives until the overlap between modes and the noise around them scramble it. In the
2025 devices the first was milliseconds and the second microseconds.

\textbf{4. Universal gates.} Clifford gates by sequences of measurements; the $T$ gate by magic
states.

\textbf{5. Measurement.} The quantum dot and its capacitance, in a single shot. In 2025 the $Z$
readout was wrong 0.5\% of the time and the $X$ readout 16\%.

\textbf{6 and 7. Moving qubits.} The same problem as transmons: everything lives in a
dilution refrigerator and talks through microwaves.

\section{No winner yet}

\begin{center}
\footnotesize
\renewcommand{\arraystretch}{1.25}
\begin{tabular}{>{\raggedright\arraybackslash}p{22mm}*{6}{>{\raggedright\arraybackslash}p{16.5mm}}}
\hline
 & \textbf{1.} characterized, scalable & \textbf{2.} initial state & \textbf{3.} coherence per gate & \textbf{4.} gates & \textbf{5.} readout & \textbf{6--7.} moving qubits \\
\hline
super\-con\-duct\-ing & fabricated, wiring & needs mK & $\sim10^4$ & neighbors only & $\sim$99\% & hard \\
trapped ions & identical; short chains & optical pumping & $\sim10^4$ & all-to-all, slow & $>$99.9\% & natural \\
neutral atoms & identical; thousands & optical pumping & good & Rydberg, slower & good & possible \\
photons & sources, losses & easy & excellent & proba\-bil\-is\-tic & detectors & natural \\
semicond. spins & non-uniform & needs mK & improving & neighbors & improving & hard \\
topological & not shown & by mea\-sure\-ment & -- & Clifford only & -- & hard \\
\hline
\end{tabular}
\end{center}

Every platform trades one criterion against another. The transmon is easy to make and hard to
wire; the ion is perfect and slow; the atom is plentiful and loosely coupled; the photon travels
well and does not interact. Charts that score each platform from 0 to 5 circulate widely, and they
are opinions dressed as data. The table above says something more modest and more durable: none
of them has yet met all five criteria at the scale that error correction needs, which is the
subject of Part~\ref{part:qec}.
